\section{Limitations, corrections, and the next falsifiable step}
\subsection{Completeness and certification obligations}
A generalized replacement must cover the input system, not just this norm
chart. The remaining bridge includes original summation-chain auxiliary
fibers, exceptional maps, $z=0$, rational lifts versus conjugate or
ghost norm configurations, and all branch multiplicities. Each branch
needs an exact original-system validator and a coverage argument. Testing
a supplied $H$ or rejecting sampled relaxed norms cannot close the
coefficient-search problem: the recorded divisor relaxation still has
$2^{73}$ candidates at fixed parameters \ev{E1}.

Theorem~\ref{thm:support} is valid over every finite field.
Theorem~\ref{thm:chart} is explicitly characteristic two. Koblitz curves are
a structured binary subclass, not a third field characteristic. Applicable
binary formulas still require a field representation, valid curve parameters,
and the theorem's target hypotheses. Prime curves require an odd-characteristic
chart and an appropriate domain representation. Large-prime supplied-quintic
validation is useful infrastructure, but does not fill that gap.
No statement here establishes an efficient universal algorithm at arbitrary
sizes.

All theorems have readable proofs in this manuscript, but no checked Lean
formalization is supplied. The local commands \texttt{lean} and
\texttt{conductor} are unavailable in the manuscript environment; that is a
specific tooling limitation, not a proof or a reservation. The central
formalization obligation is the coefficient identity together with its
finite-field square-root and Artin--Schreier hypotheses. The separate existing
Lean proof in the earlier decomposition manuscript is not a proof of these
theorems. Appendix~\ref{app:artifacts} maps the outstanding obligations.

The repository's current \texttt{identity.certificate/v1} checker evaluates
unconditional identities modulo $2^{61}-1$. It cannot faithfully certify the
characteristic-two identities or conditional implications here by merely
changing a label. We therefore attach source-hashed native evidence pointers
and the written proofs, with this certificate-format gap stated explicitly.
No such certificate or formal verification is claimed.

\subsection{Preserved erratum}
A prior transfer-only proof omitted a term. The correction \ev{E6} is
\[
 \begin{split}
 V=e_5+e_1e_4
   &=(h^2+h)e_4\\
   &\quad+k(zs+z^2r+ky_T+z^3+hz y_T).
 \end{split}
\]
Thus $k=0$ does not generally imply $V=0$.
Over $\F_4$ with $\omega^2=\omega+1$, take $a=0,b=1$, target
$(z,y_T)=(1,0)$, and $h=\omega,k=r=s=0$. Literal division gives
$H=X^5+X+1$, so $e_1=0$ and $e_4=e_5=1$, hence $V=1$.
This counterexample invalidates the omitted-term formula.
Generic support code computes the coefficients directly and does not use
that simplification. Historical files remain preserved; the correction
must travel with any use of the old transfer argument.

\subsection{Performance gate and comparison still owed}
The program requires twenty preregistered hard matching instances, completed
all-roots outputs on every method, a fastest baseline taking at least one
second, and at least three cold repetitions. Genuine source-identified
F4/F5/F6 interfaces must return the same original roots. The immutable floor
is at least a twofold geometric-mean end-to-end gain, no more than 1.10 worst
slowdown, A/A variation at most 10\%, qualified isolation, and the requested
field coverage \ev{E7}. The higher requested speed targets are research
objectives, not achieved results. Nothing in this manuscript passes the floor.

Comparisons still owed include strong available systems such as Magma,
msolve, FGb, Singular, and OpenF4, subject to field and contract support.
The supplied-polynomial stage also owes a matched comparison against
established factorization implementations such as FLINT, NTL, or Magma.
Baseline tuning, conversion, setup, solving, extraction, verification, and
serialization must all be charged consistently. Unsupported fields are
reported, not silently replaced with easier ideals. Native SAT is not
itself an F4/F5/F6 baseline.

\subsection{Next queue}
The failed prediction leads to a changed experiment, not an identical rerun.
The first queued variant is sorted-source norm-primary search:
add an exact comparator network for the canonical source encoding, prove its
unique extensions, enumerate representatives, and restore every distinct
permutation. Before a hard run, require exact agreement with all eight
symmetry-oracle sets \ev{E4}, including zero and repeated sources.
Freeze the comparator source, input digest, native source manifest, budgets,
and decision mask. Predict a verified witness or representative exhaustion
on the same actual fixture within a bounded inner search; on failure,
retain UNKNOWN and move to a structural variant. At most three experiments
may be scheduled in a bounded session, each at one core and 4 GiB.

The second route is conditional whole-cell pruning using the support
nilpotence equations, rather than enumerating all relaxed divisors.
A proposed cell certificate must imply rejection of every point in that
cell; rejection of random sample norms is insufficient. Its small oracle
must compare complete cell membership and cover repeated factors.
The third route develops an odd-characteristic norm parametrization with
a separate prime-domain contract. It must preserve odd cross terms and
prove exceptional branches before importing any binary optimization.

This queue separates model preservation, useful structural pruning, and
end-to-end performance. The present evidence improves representation and
rejects a bounded search prediction. It leaves both the complete generalized
algorithm and its qualified speed advantage unresolved.
